% from ega2-1-fr.tex
\oldpage[II]{138}

où $\mathcal{G}$ et $\operatorname{Coker}v$ sont des
$\mathcal{O}_Y$-Modules cohérents
(\hyperref[0.5.3.4-fr]{0, 5.3.4} et
\hyperref[0.5.3.5-fr]{5.3.5})~; en vertu de la suite exacte de cohomologie,
il suffira de montrer que
$\mathrm{H}^1(Y,\mathcal{G})
=\mathrm{H}^1(Y,\operatorname{Coker}v)=0$ pour en déduire
$\mathrm{H}^1(Y,\mathcal{F}^m)
=(\mathrm{H}^1(Y,\mathcal{F}))^m=0$, et par suite
$\mathrm{H}^1(Y,\mathcal{F})=0$. Or la restriction $v|V$ est un isomorphisme,
donc le support de $\operatorname{Coker}v$ est distinct de $Y$, d'où
$\mathrm{H}^1(Y,\operatorname{Coker}v)=0$ par l'hypothèse du début. D'autre
part, $\mathcal{G}$ est un $\mathcal{B}$-Module cohérent
(\hyperref[I.9.6.4-fr]{I, 9.6.4})~; comme $X$ est affine au-dessus de $Y$, il
existe un $\mathcal{O}_X$-Module quasi-cohérent $\mathcal{K}$ tel que
$\mathcal{G}$ soit isomorphe à $f_*(\mathcal{K})$
(\hyperref[II.1.4.3-fr]{1.4.3})~; comme $X$ est affine, on a
$\mathrm{H}^1(X,\mathcal{K})=0$
(\hyperref[I.5.1.9.2-fr]{I, 5.1.9.2}), donc aussi
$\mathrm{H}^1(Y,\mathcal{G})=0$ par
(\hyperref[II.5.2.3-fr]{5.2.3}), ce qui achève la démonstration du théorème
(\hyperref[II.6.7.1-fr]{6.7.1}).

\begin{corollary}[6.7.2]
\label{II.6.7.2-fr}
Soient $X$ un préschéma noethérien, $(X_i)_{1\leq i\leq n}$ un recouvrement
fini de l'espace $X$ formé d'ensembles fermés. Pour que $X$ soit affine, il
faut et il suffit que pour chaque $i$, il existe un sous-préschéma fermé de
$X$, ayant $X_i$ pour espace sous-jacent et qui soit affine.
\end{corollary}

En effet, s'il en est ainsi, soit $X'$ le schéma \emph{somme} des $X_i$~; il
est clair que $X'$ est affine, et on définit un morphisme surjectif
$f:X'\to X$ en prenant pour restriction de $f$ à $X_i$ l'injection
canonique. Tout revient à vérifier que $f$ est \emph{fini}, en vertu de
(\hyperref[II.6.7.1-fr]{6.7.1}), et cela a été vu en
(\hyperref[II.6.1.5-fr]{6.1.5}).

\begin{corollary}[6.7.3]
\label{II.6.7.3-fr}
Pour qu'un préschéma noethérien $X$ soit affine, il faut et il suffit que les
sous-préschémas fermés réduits ayant pour espaces sous-jacents les composantes
irréductibles de $X$ soient affines.
\end{corollary}

\section{Critères valuatifs}
\label{section:II.7-fr}

Nous donnons dans ce paragraphe des critères valuatifs de séparation et de
propreté pour un morphisme, c'est-à-dire des critères qui font intervenir un
schéma auxiliaire variable de la forme $\operatorname{Spec}(A)$, où $A$ est
un anneau de valuation. Moyennant des hypothèses «~noethériennes~»
convenables, on peut dans ces critères se restreindre au cas où $A$ est un
anneau de valuation \emph{discrète}. Ce sera sans doute le seul cas que nous
aurons à appliquer par la suite, et nous n'avons introduit les anneaux de
valuation quelconques, dans le cas général, que pour faire le lien avec des
développements classiques.

\subsection{Rappels sur les anneaux de valuation.}
\label{subsection:II.7.1-fr}

\begin{env}[7.1.1]
\label{II.7.1.1-fr}
Parmi les diverses propriétés équivalentes qui caractérisent les anneaux de
valuation, nous utiliserons la suivante~: un anneau $A$ est dit
\emph{anneau de valuation} si c'est un anneau intègre qui n'est pas un
corps, et si, dans l'ensemble des anneaux locaux contenus dans le corps des
fractions $K$ de $A$ et distincts de $K$, $A$ est \emph{maximal} pour la
relation de domination (\hyperref[I.8.1.1-fr]{I, 8.1.1}). Rappelons qu'un
anneau de valuation est \emph{intégralement clos}. Si $A$ est un anneau de
valuation, $A_{\mathfrak{p}}$ est aussi un anneau de valuation pour tout
idéal premier $\mathfrak{p}\neq0$ de $A$.
\end{env}

\begin{env}[7.1.2]
\label{II.7.1.2-fr}
Soient $K$ un corps, $A$ un sous-anneau local de $K$ qui n'est pas un corps~;

